\section{Introduction}\label{sec:introduction}

Uniform property $\Gamma$ asks for central projections that divide every
trace while preserving its values on prescribed coefficients. This form of
central divisibility is central to the regularity theory of nuclear
$C^*$-algebras. It was introduced by Castillejos, Evington, Tikuisis, White
and Winter \cite[Definition~2.1]{CETWW}, and developed systematically by
Castillejos, Evington, Tikuisis and White \cite{CETW}. For simple separable
unital nuclear algebras with traces, strict comparison together with uniform
property $\Gamma$ is equivalent to absorption of the Jiang--Su algebra
\cite[Theorem~5.6]{CETW}. We give three independent constructions of this central divisibility.
The first derives uniform property $\Gamma$ from real rank zero of a
completion ultrapower, answering a question of Schafhauser, Tikuisis and
White. The second derives Jiang--Su absorption from strict comparison
in the unital setting, completing the unital Toms--Winter equivalence.
The third treats stably projectionless algebras with bounded densely
finite traces and a compact normalized trace base. The common task is
to make projections central while preserving every coefficient trace;
the hypotheses and the averaging arguments differ in the three routes.

The underlying finite-factor notion goes back to Murray and von Neumann
\cite{MurrayVonNeumann1943}. Dixmier developed its central projection
formulation \cite{Dixmier1969}, while McDuff studied the stronger central
matrix structure associated with absorption of the hyperfinite factor
\cite{McDuff1970}. Connes' classification of injective factors
\cite{Connes} supplies the pointwise von Neumann regularity used in our
nonunital construction. The uniform $C^*$-algebraic question asks for a
single family of projections that works simultaneously across the trace
space and retains every coefficient identity. Pointwise tracial
approximation must therefore be combined with a uniform construction.

The first construction starts with real rank zero in a uniform tracial
ultrapower. A unital $C^*$-algebra has \emph{real rank zero} when its
selfadjoint elements can be approximated in norm by selfadjoint elements
with finite spectrum \cite{BP91}. We first describe precisely
the resulting trace identities. For a unital algebra $A$ with nonempty
tracial-state space $T(A)$, put
\[
 \|a\|_{2,T(A)}=\sup_{\tau\in T(A)}\tau(a^*a)^{1/2}.
\]
The uniform tracial completion belongs to the framework initiated by
Ozawa \cite{Ozawa13} and developed in \cite{CCEGSTW}.
The algebra $B=\overline A^{T(A)}$ consists of bounded
uniformly $2$-Cauchy sequences modulo sequences tending to zero in this
seminorm. Each $\tau\in T(A)$ extends to a designated trace $\bar\tau$ on $B$.
For a fixed free ultrafilter $\omega$ on $\N$, the uniform tracial ultrapower
$B^\omega$ is the quotient of bounded sequences in $B$ by those whose uniform
$2$-norm tends to zero along $\omega$. The family $\Lambda_\omega$ consists
of exactly the traces
\[
 \lambda([(b_j)])=\lim_{j\to\omega}\bar\tau_j(b_j),
 \qquad \tau_j\in T(A).
\]
These constructions and their quotient norms are specified in
Section~\ref{r:sec:preliminaries}.

\begin{theorem}[Real rank zero and uniform property $\Gamma$]
\label{r:thm:main}
Let $A$ be a unital, separable, simple, infinite-dimensional, nuclear,
stably finite $C^*$-algebra with $T(A)\ne\varnothing$, and let
$B=\overline A^{T(A)}$. Fix any free ultrafilter $\omega$ on $\N$.
If $B^\omega$ has real rank zero in its quotient $C^*$-norm, then there is
a projection $p\in B^\omega\cap B'$ such that
\[
 \lambda(px)=\tfrac12\lambda(x)
 \qquad(x\in B,\ \lambda\in\Lambda_\omega).
\]
In particular, $A$ has uniform property $\Gamma$.
\end{theorem}

This gives a positive resolution of Problem~XXI of Schafhauser, Tikuisis
and White \cite[Section~6]{STW}. Their Definition~16 uses precisely this
completion and halving formulation. In the usual formulation, uniform
property $\Gamma$ requires, for every $m\ge1$, orthogonal projections
$p_1,\ldots,p_m\in A^\omega\cap A'$ summing to one such that
\[
 \lambda(ap_i)=\frac1m\lambda(a)
 \quad(a\in A,\ \lambda\in T_\omega(A)).
\]
Here $A^\omega$ is the uniform tracial ultrapower formed directly from $A$,
and $T_\omega(A)$ is its family of varying limit traces. The completion
formulation and the all-$m$ formulation agree by
\cite[Propositions~5.7, 5.20 and 5.23]{CCEGSTW}; we record the interfaces
beside the definitions in Section~\ref{r:sec:preliminaries}. The coefficient $a$ is essential: equality
of scalar traces of the $p_i$ alone does not give uniform property $\Gamma$.

The additional centrality in Theorem~\ref{r:thm:main} is the main issue.
Vaccaro derives tracial almost divisibility from real rank zero of the
tracial ultrapower: hereditary subalgebras admit tracially large matrix
order-zero maps, without a centrality requirement
\cite[Proposition~2.2]{Vac26}. Combined with tracially locally finite nuclear
dimension, this yields uniform McDuffness and hence uniform property
$\Gamma$ \cite[Theorem~2.3]{Vac26}. Kessinger and Toms obtain uniform
property $\Gamma$ from subquadratic growth in recursive-subhomogeneous
inductive limits and from a tracial local approximation condition
\cite[Theorems~1.2 and 7.11]{KT26}. Their results exclude nonzero
finite-dimensional representations; the local theorem also assumes
nonempty traces. Their argument supplies both ultrapower real rank zero
and tracially locally finite nuclear dimension, then uses this divisibility
route. We instead construct central coefficient splittings using
nuclear averaging under the hypotheses of Theorem~\ref{r:thm:main}.
Nuclearity alone does not ensure uniform property $\Gamma$: Toms constructs
a simple separable unital AH algebra without it \cite[Theorem~1]{Toms26}.

There are also implications in the other direction. For simple separable
unital infinite-dimensional nuclear algebras with traces, Fu proves that
uniform property $\Gamma$ implies the real-rank and divisibility
conditions of \cite[Theorems~1.1 and 10.6]{Fu26}.
Uniform property $\Gamma$ also passes to the tracial completion and
supplies complemented partitions of unity
\cite[Proposition~5.20 and Theorem~1.4]{CCEGSTW}.
For factorial tracially complete algebras, complemented partitions of unity
imply real rank zero of their tracial ultrapowers
\cite[Proposition~7.2]{CCEGSTW}; Evington and Tikuisis strengthen this to
real rank zero of the tracially complete algebra itself
\cite[Theorem~A]{EvingtonTikuisis26}. These results assume the uniform
partition structure from which projections can be assembled.
Theorem~\ref{r:thm:main} constructs central splittings from the ultrapower
real-rank-zero assumption.

The same hypotheses also identify all traces on the completion.

\begin{corollary}[Traces on the completion]\label{r:cor:completion-traces}
Under the hypotheses of Theorem~\ref{r:thm:main}, including real rank zero
of $B^\omega$ in its quotient $C^*$-norm, every tracial state of $B$ is
the designated extension of a unique trace in $T(A)$. Thus restriction
is an affine homeomorphism $T(B)\to T(A)$, with inverse
$\tau\mapsto\bar\tau$. Moreover,
\[
 |\sigma(b)|\le\|b\|_{2,T(A)}
 \qquad(b\in B,\ \sigma\in T(B)),
\]
so every tracial state of $B$ is continuous for the designated uniform
$2$-norm.
\end{corollary}

The proof appears at the end of Section~\ref{r:sec:completion}.

The construction is useful outside an ultrapower. A projection is called
\emph{full} when the closed two-sided ideal it generates is the entire
ambient algebra.

\begin{theorem}[Finite-set central splitting]\label{r:thm:finite}
Let $A\subseteq D$ be a unital inclusion of $C^*$-algebras, with $A$ nuclear,
$D$ of real rank zero, and $T(D)\ne\varnothing$. Suppose that for every
$\delta>0$ there is a full projection $g\in D$ with
$\sup_{\tau\in T(D)}\tau(g)<\delta$.
For every finite $K\subseteq\mathcal U(A)$ and $\varepsilon>0$ there is a
projection $P\in D$ such that, for every $x\in K$,
\[
 \sup_{\tau\in T(D)}\|[P,x]\|_{2,\tau}<\varepsilon,
 \qquad
 \sup_{\tau\in T(D)}|\tau(Px)-\tfrac12\tau(x)|<\varepsilon.
\]
\end{theorem}

The ambient algebra in this result need not be nuclear or simple. Its
tracial states need not be faithful. In the application $D=B^\omega$,
simplicity and infinite dimensionality of $A$ provide the full projections
of small trace uniformly over \emph{all} abstract tracial states of $D$.
The last diagonal argument then uses the specified family
$\Lambda_\omega$ to obtain a single projection commuting with all of $B$.

\subsection{Strict comparison and Jiang--Su absorption}

The next two results obtain uniform property $\Gamma$ directly from
strict comparison. Each supplies its own projections and is independent
of the real-rank-zero hypothesis in Theorem~\ref{r:thm:main}.
For positive $a,b\in A\otimes\K$, Cuntz subequivalence $a\precsim b$ means
that $x_n^*bx_n\to a$ in norm for a sequence $x_n\in A\otimes\K$.
Here $\K$ denotes the compact operators, tensor products are spatial,
and each trace is extended using the unnormalized matrix trace. Set
\[
 d_\tau(a)=\lim_{n\to\infty}(\tau\otimes\Tr)(a^{1/n})\in[0,\infty].
\]
For simple unital algebras we use the comparison convention
\[
 b\ne0,\quad d_\tau(a)<d_\tau(b)\ (\tau\in T(A))
 \quad\Longrightarrow\quad a\precsim b.
\]
If $T(A)$ is empty, this convention requires comparison with every nonzero
positive target; it is more than the absence of traces.

The Jiang--Su algebra $\Z$ was introduced by Jiang and Su \cite{JS};
$\Z$-stability means $A\cong A\otimes\Z$. For simple, separable,
unital, infinite-dimensional nuclear algebras, the Toms--Winter
conjecture asks for the equivalence \cite[Conjecture~9.3]{WZNuc}
\[
 \text{finite nuclear dimension}
 \quad\Longleftrightarrow\quad
 \mathcal Z\text{-stability}
 \quad\Longleftrightarrow\quad
 \text{strict comparison}.
\]
Finite nuclear dimension requires completely positive approximations
through finite-dimensional algebras with a fixed finite number of
outgoing pieces that preserve orthogonality of positive elements.
The implications from finite nuclear dimension to $\Z$-stability,
from $\Z$-stability to strict comparison, and from $\Z$-stability to
finite nuclear dimension are known \cite{WinterNuc,Rordam,CETWW}.
Theorem~\ref{u:thm:main} proves the remaining comparison-to-absorption
implication, giving a positive resolution in the unital setting.
The empty-trace case is included with the comparison convention stated
above; its purely infinite absorption argument is given in
Section~\ref{u:sec:gamma}.

\begin{theorem}[The direct unital comparison route]\label{u:thm:main}
Let $A$ be a simple, separable, unital, infinite-dimensional nuclear
$C^*$-algebra with strict comparison in the preceding sense.
Then $A\cong A\otimes\Z$. When $T(A)\ne\varnothing$, the proof constructs
uniform property $\Gamma$ directly.
\end{theorem}

In the tracial case, Matui and Sato established the absorption implication
for finitely many extremal traces using comparison and property~(SI)
\cite[Theorem~1.1]{MS}.
The compact finite-dimensional boundary case was obtained independently
by Kirchberg and R\o rdam \cite[Corollary~7.9]{KR14}, Sato
\cite[Corollary~1.2]{Sato12}, and Toms, White and Winter
\cite[Corollary~4.7]{TWW15}. Later results reach other trace spaces: Lin
proves the implication under his condition~(C) together with stable rank
one, and removes the stable-rank assumption in its countable-boundary case
\cite[Theorem~5.6 and Corollary~5.7]{Lin24}. Condition~(C) decomposes the
extreme boundary into countably many finite-dimensional compact pieces
with additional compatibility requirements on their convex faces
\cite[Definition~2.13]{Lin24}. Evington and Schafhauser subsequently
obtained uniform property $\Gamma$ from fibrewise property $\Gamma$
under the compact finite-dimensional boundary hypothesis without assuming
nuclearity \cite[Theorem~1.1]{EvingtonSchafhauser26}. The uniform divisibility criterion
of \cite[Theorem~5.6]{CETW} explains why constructing central projections
is the decisive task for the comparison-to-absorption implication.

For the nonunital result, write $T_{\mathrm{lsc}}(A)$ for the cone of densely
finite lower-semicontinuous traces and $T_1(A)$ for its bounded traces of
norm one. The next theorem has a distinct trace hypothesis.

\begin{theorem}[The direct stably projectionless route]\label{n:thm:main}
Let $A$ be a nonzero simple, separable, nuclear $C^*$-algebra such that
$A\otimes\K$ has no nonzero projections. Suppose that every trace in
$T_{\mathrm{lsc}}(A)$ is bounded on the positive unit ball and that
$T_1(A)$ is nonempty and weak-$*$ compact. Assume that, for
$a,b\in(A\otimes\K)_+$ with $b\ne0$,
\[
 d_\tau(a)<1\quad
 \text{for every }\tau\in T_{\mathrm{lsc}}(A)
 \text{ with }d_\tau(b)=1
 \quad\Longrightarrow\quad a\precsim b.
\]
Then $A$ has uniform property $\Gamma$ and $A\cong A\otimes\Z$.
\end{theorem}

Boundedness of the entire densely finite trace cone and compactness of its
normalized base are separate assumptions. The comparison tests are exactly
those trace rays with positive finite target rank; ranks in the stabilization
may be infinite. The construction leads through the nonunital uniform
McDuffness and absorption criteria of Castillejos, Evington, Tikuisis and
White \cite{CETW} and Castillejos, Li and Szab\'o \cite{CLS}. Jacelon's work
on nonunital algebras with finite-dimensional trace boundary and continuous
rank functions is an earlier part of this development \cite{Jacelon}.

\subsection{The three constructions and their interfaces}

The common problem is to turn projections that split traces into projections
that are central and split \emph{coefficient} traces. Figure~\ref{fig:routes}
shows the separate routes. Each is proved in full, and the first uses no
comparison assumption.

\begin{figure}[htbp]
\centering
\begin{tikzpicture}[node distance=5mm and 5mm, every node/.style={font=\small},
 box/.style={draw,rounded corners=2pt,align=center,text width=42mm,
 minimum height=13mm,inner sep=4pt}, >={Stealth[length=2mm]}]
\node[box] (r1) {Real rank zero\\and nuclear averaging};
\node[box,right=of r1] (r2) {Weighted blocks\\and orthogonal packing};
\node[box,right=of r2] (r3) {Halving in $B^\omega$\\against all $x\in B$};
\node[box,below=of r1] (u1) {Unital comparison\\and rank rounding};
\node[box,right=of u1] (u2) {Centered symmetries\\and a Gaussian limit};
\node[box,right=of u2] (u3) {Uniform $\Gamma$\\and absorption};
\node[box,below=of u1] (n1) {Nonunital comparison\\with bounded traces\\and a compact base};
\node[box,right=of n1] (n2) {Cost selection\\and bounded truncation};
\node[box,right=of n2] (n3) {Uniform $\Gamma$\\and nonunital absorption};
\draw[->] (r1)--(r2); \draw[->] (r2)--(r3);
\draw[->] (u1)--(u2); \draw[->] (u2)--(u3);
\draw[->] (n1)--(n2); \draw[->] (n2)--(n3);
\end{tikzpicture}

\caption{Three routes to central coefficient-tested trace splitting. The
upper route uses real rank zero of the completion ultrapower. The lower
routes use comparison and supply their own projections in $A^\omega$.
Absorption is a consequence of the two comparison routes through the
published criteria verified in the text.}
\label{fig:routes}
\end{figure}

For the real-rank-zero route, Haagerup's convex virtual diagonal
\cite[Theorem~3.1]{Haa83} provides nuclear averaging maps chosen before
any input projection. Finite scalar compressions turn their outputs into
weighted matrix blocks. A small proportion of a base projection can be
sampled inside each block without making that subprojection central:
scalar matrix compressions retain the coefficient traces and squared
commutators. Orthogonal insertion then packs two labelled families with matching
target traces. The estimates use total weighted mass and energy, so they remain
independent of the number and sizes of the blocks. A bounded diagonal
argument yields the desired projection for the original ultrafilter, and
uniform $2$-density extends its identities from $A$ to its completion.

For the unital comparison route, a full positive element of small rank
receives an exact trace-preserving compression of a given positive element.
Integer spectral cuts approximate that trace by a finite-matrix rank, as
proved in Section~\ref{u:sec:ranks}; this elementary step needs neither
nuclearity nor comparison. Comparison then produces projections with exact
spectral support in $A^\omega$. Scalar compressions and matrix-unit transport
construct symmetries centered against any prescribed separable coefficient
algebra. The convex order-zero approximations of Hirshberg, Kirchberg and
White \cite[Theorem~1.4]{HKW} give a square-root-weighted sum of these
symmetries with fourth moment at most three. Its norm may grow with the
number of summands. Clipping at a fixed radius supplies bounded relative
moment estimates; successively adjoining commuting variables gives normalized sums with a
common Gaussian limit in coefficient traces. Continuous approximations to
equal-measure Gaussian intervals then yield the required central projections.

For the stably projectionless route, spectral allocation supplies small
projections with a one-sided cost bound. A single cost records both overlap
with earlier projections and error in a finite-dimensional model. A
block-matrix functional-calculus estimate passes aggregate squared
commutators to a truncated sum without a factor for the number of models.
The resulting small splittings have trace and projection errors of order
$\lambda^{3/2}$. Reindexing and multiplication amplify them to exact
splittings. The final passage distinguishes the uniform tracial ultrapower
$A^\omega$, the norm ultrapower $A_\omega$, and the annihilator quotient
$F_\omega(A)$, and checks generalized limit traces explicitly.

Sections~\ref{r:sec:preliminaries}--\ref{r:sec:completion} prove the
real-rank-zero route, including the finite-set theorem.
Sections~\ref{u:sec:preliminaries}--\ref{u:sec:gamma} give the direct unital
route and handle its empty-trace case. Sections~\ref{n:sec:framework}--
\ref{n:sec:absorption} give the bounded-trace nonunital route.
The elementary small-support lemma in Section~\ref{r:sec:projections}
serves both unital routes; Section~\ref{u:sec:ranks} proves the
trace-transfer and integer-rounding argument at its point of use.
